% figures/methodology/tab_workflow.tex — the detect-translate-evaluate workflow
% Requires \usepackage{booktabs}. Shows phase, method, input, output.
\begin{table}[!ht]
    \centering
    \small
    \renewcommand{\arraystretch}{1.3}
    \caption[The detect--translate--evaluate workflow.]{The thesis pipeline. Each
    phase takes the output of the previous one as its input, so that the analysis
    constrains the design and the design is what the evaluation tests.}
    \label{tab:workflow}
    \begin{tabular}{@{}>{\raggedright\arraybackslash}p{2.3cm}>{\raggedright\arraybackslash}p{3.5cm}>{\raggedright\arraybackslash}p{3.2cm}>{\raggedright\arraybackslash}p{3.2cm}@{}}
        \toprule
        \textbf{Phase} & \textbf{Method} & \textbf{Input} & \textbf{Output} \\
        \midrule
        \textbf{Detect}
            & PCA, autoencoder, anomaly \& early-warning analysis
            & RAPID observations; Copernicus reanalysis
            & Interpretable latent structure of the system \\
        \textbf{Translate}
            & Design grammar applied to recovered features; browser-based build
            & Latent structure from \emph{detect}
            & Three representational conditions (static, dynamic, experiential) \\
        \textbf{Evaluate}
            & Two-group study; survey-based evaluation; exploratory analysis
            & Representations from \emph{translate}
            & Evidence on how non-specialists understand the system \\
        \bottomrule
    \end{tabular}
\end{table}